\documentclass[10pt,a4paper]{amsart}
\usepackage[margin=19mm]{geometry}
\usepackage{amsmath,amssymb,mathtools}
\usepackage[scr=rsfs,cal=euler]{mathalfa}
\usepackage{xfrac}
\usepackage[unicode,hidelinks,bookmarks=false]{hyperref}
\newcommand{\ie}{i.e.\ }
\newcommand{\cf}{cf.\ }
\newcommand{\resp}{respectively\ }
\newcommand{\Subsection}{\subsection}
\providecommand{\nobreakdash}{\nobreak\hbox{-}}
\setlength{\emergencystretch}{2em}
\multlinegap0pt
\usepackage{mathtools}
\usepackage{enumerate}
\usepackage{amssymb}
\usepackage{tikz}
\newtheorem{theorem}{Theorem}
\title{This is the title}
\author{K. Mahesh Krishna}
\date{Source: arXiv:2506.18913v2 (CC BY 4.0)}
\begin{document}
\maketitle

\noindent\emph{Source and licence.} This is the title. Version arXiv:2506.18913v2 (CC BY 4.0), distributed by arXiv under the Creative Commons Attribution 4.0 International licence, http://creativecommons.org/licenses/by/4.0/, as recorded in the arXiv licence metadata for this exact version and shown on the abstract page https://arxiv.org/abs/2506.18913v2. Full licence notice: \texttt{licenses/arxiv-2506.18913-CC-BY-4.0.txt}.\par\smallskip

\noindent\textit{Editorial scope.} Counted unit: the article's theorem NGJ (source0.tex lines 385--396). The article's complete original proof of it is delivered with it (source0.tex lines 397--457, one \texttt{proof} environment of 61 lines). No mathematics is written, repaired or re-derived: the statement and the proof are the article's own lines, in the article's own notation, and no retained line is substituted or re-flowed.  No further source-local context is needed: every cross-reference of every retained line resolves inside the two delivered spans.  Nothing was omitted from any retained span, so the package carries no omission record.  No retained line contains a citation, so the wrapper carries no bibliography and no bibliography entry.  No retained line contains a figure environment, an image-inclusion command or drawing code, so no image is delivered and the package declares no asset dependency.  Editorial formatting only, in addition: an AMS \texttt{amsart} preamble that restates the article's own declarations and macros used by the retained lines, namely the environments \texttt{theorem} on separate counters, together with the standard packages those lines need.  The retained environments print in this document as Theorem 1 in order of appearance, whereas the article prints the same items with the numbering of its own full document.  The complete native source of the article as distributed in the arXiv e-print for this exact version is bundled unchanged as \texttt{sources/180374/source0.tex}.
\begin{theorem}\label{NGJ}
	(\textbf{Non-Archimedean  Ghobber-Jaming  Uncertainty Principle})
	Let $\mathcal{X}$  be a finite dimensional non-Archimedean Banach  space over $\mathbb{K}$.	Let $(\{f_j\}_{j=1}^n, \{\tau_j\}_{j=1}^n)$	and $(\{g_k\}_{k=1}^n, \{\omega_k\}_{k=1}^n)$ be orthonormal bases  for $\mathcal{X}$.  If $M,N\subseteq \{1, \dots, n\}$ are such that 
	\begin{align*}
		\displaystyle \max_{j \in M, k \in N}|g_k(\tau_j)|<1,
	\end{align*}
	then for all $x \in \mathcal{X}$, 	
	\begin{align}\label{NGH}
			\|x\|\leq  \left(\frac{1}{1-\displaystyle \max_{j \in M, k \in N}|g_k(\tau_j)|}\right)\max\left\{\displaystyle \max_{j \in M^c}|f_j(x) |, \displaystyle \max_{k \in N^c}|g_k(x)|\right\}.
	\end{align}
	In particular, if $x$ is supported on $M$ in the expansion using basis $(\{f_j\}_{j=1}^n, \{\tau_j\}_{j=1}^n)$	 and $x$ is supported on $N$ in the expansion using basis $(\{g_k\}_{k=1}^n, \{\omega_k\}_{k=1}^n)$, then $x=0$. 
\end{theorem}

\begin{proof}
Given $S\subseteq \{1, \dots, n\}$, define 

\begin{align*}
	&P_Sx\coloneqq \sum_{j\in S}f_j(x)\tau_j, \quad \forall x \in \mathcal{X}, \\
	&\|x\|_{S, f}\coloneqq \displaystyle \max_{j \in S}|f_j(x)|, \quad \forall x \in \mathcal{X},\\
	&\|x\|_{S, g}\coloneqq \displaystyle \max_{j \in S}|g_j(x)|, \quad \forall x \in \mathcal{X}.
\end{align*}
Then 
\begin{align*}
	\|P_Sx\|=\left\|\sum_{j\in S}f_j(x)\tau_j\right\|=\displaystyle \max_{j \in S}|f_j(x)|=	\|x\|_{S, f}, \quad \forall x \in \mathcal{X}.
\end{align*}
Define $V:\mathcal{X}\ni x \mapsto \sum_{k=1}^{n}g_k(x)\tau_k \in \mathcal{X}$. Then $V$ is an invertible isometry. Using $V$ we get 
\begin{align*}
	\|P_SVx\|&=\left\|\sum_{j\in S}f_j(Vx)\tau_j\right\|=\left\|\sum_{j\in S}f_j\left(\sum_{k=1}^{n}g_k(x)\tau_k\right)\tau_j\right\|=\left\|\sum_{j\in S}\sum_{k=1}^{n}g_k(x)f_j(\tau_k)\tau_j\right\|\\
	&=\left\|\sum_{j\in S}g_j(x)\tau_j\right\|=\displaystyle \max_{j \in S}|g_j(x)|=	\|x\|_{S, g}, \quad \forall x \in \mathcal{X}.
\end{align*}
For $x \in \mathcal{X}$, we  now find 
\begin{align*}
	\|P_NVP_Mx\|&=\left\|\sum_{k\in N}f_k(VP_Mx)\tau_k\right\|=\max_{k \in N}|f_k(VP_Mx)|\\
	&=\max_{k \in N}\left|(f_kV) \left(\sum_{j\in M}f_j(x)\tau_j\right)\right|=\max_{k \in N}\left| \sum_{j\in M}f_j(x)f_k(V\tau_j)\right|\\
	&=\max_{k \in N}\left| \sum_{j\in M}f_j(x)f_k\left(\sum_{r=1}^{n}g_r(\tau_j)\tau_r\right)\right|
	=\max_{k \in N}\left| \sum_{j\in M}f_j(x)\sum_{r=1}^{n}g_r(\tau_j)f_k(\tau_r)\right|\\
	&=\max_{k \in N}\left| \sum_{j\in M}f_j(x)g_k(\tau_j)\right|\leq  \max_{k \in N}\max_{j \in M}|f_j(x)g_k(\tau_j)|\\
	&\leq \left(\max_{j \in M, k\in N}|g_k(\tau_j)|\right)\max_{j \in M}|f_j(x)|\leq \left(\max_{j \in M, k\in N}|g_k(\tau_j)|\right)\|x\|.
\end{align*}
So we have 
\begin{align}\label{FPVP}
	\|P_NVP_M\|\leq \displaystyle \max_{j \in M, k\in N}|g_k(\tau_j)|.
\end{align}
Now let $y \in \mathcal{X}$ be such that $\{j \in \{1, \dots, n\}: f_j(y)\neq 0\} \subseteq M.$
Then 
\begin{align*}
	\|P_NVy\|=\|P_NVP_My\|\leq \|P_NVP_M\|\|y\|
\end{align*}
 and 
\begin{align*}
	\|y\|_{N^c, g}&=\|P_{N^c}Vy\|=\|Vy-P_NVy\|\\
	&\geq \|Vy\|-\|P_NVy\|=\|y\|-\|P_NVy\| \\
	&\geq \|y\|-\|P_NVP_M\|\|y\|.
\end{align*}
Therefore 
\begin{align}\label{INP}
	\|y\|_{N^c, g}\geq (1-\|P_NVP_M\|)\|y\|.
\end{align}
Let $x \in \mathcal{X}$.   Note that $P_Mx$ satisfies $	\{j \in \{1, \dots, n\}: f_j(P_Mx)\neq 0\} \subseteq M.$ Now using (\ref{FPVP}) and (\ref{INP}) we get 

\begin{align*}
	\|x\|&=\|P_{M^c}x+P_Mx\|\leq \max\{\|P_{M^c}x\|, \|P_Mx\|\}\\
	&\leq \max \left\{\|P_{M^c}x\|, \frac{1}{1-\|P_NVP_M\|}	\|P_Mx\|_{N^c, g}\right\} \\
	&= \max \left\{\|P_{M^c}x\|, \frac{1}{1-\|P_NVP_M\|}	\|P_{N^c}VP_Mx\|\right\} \\
	&=\max \left\{\|P_{M^c}x\|, \frac{1}{1-\|P_NVP_M\|}	\|P_{N^c}V(x-P_{M^c}x)\|\right\} \\
	&\leq \max \left\{\|P_{M^c}x\|, \frac{1}{1-\|P_NVP_M\|}	\|P_{N^c}Vx\|, \frac{1}{1-\|P_NVP_M\|}	\|P_{N^c}VP_{M^c}x\|\right\} \\
	&\leq \max \left\{\|P_{M^c}x\|, \frac{1}{1-\|P_NVP_M\|}	\|P_{N^c}Vx\|, \frac{1}{1-\|P_NVP_M\|}	\|P_{M^c}x\|\right\} \\
	&=\max \left\{\frac{1}{1-\|P_NVP_M\|}	\|P_{M^c}x\|, \frac{1}{1-\|P_NVP_M\|}	\|P_{N^c}Vx\|\right\} \\
	&=\left(\frac{1}{1-\|P_NVP_M\|}\right)\max \left\{\|P_{M^c}x\|, \|P_{N^c}Vx\|\right\} \\
	&=\left(\frac{1}{1-\|P_NVP_M\|}\right)\max \left\{\|x\|_{M^c, f}, \|x\|_{N^c, g}\right\} \\
	&= \left(\frac{1}{1-\|P_NVP_M\|}\right)\max\left\{\displaystyle \max_{j \in M^c}|f_j(x)|, \displaystyle \max_{k \in N^c}|g_k(x)|\right\}\\
	&\leq  \left(\frac{1}{1-\displaystyle \max_{j \in M, k \in N}|g_k(\tau_j)|}\right)\max\left\{\displaystyle \max_{j \in M^c}|f_j(x) |, \displaystyle \max_{k \in N^c}|g_k(x)|\right\}.
\end{align*}
\end{proof}

\end{document}
